% \documentclass[letterpaper, 10pt]{article}
\documentclass[11pt, a4paper]{article}
% \documentclass{letter}
% \documentclass[runningheads]{llncs}

\usepackage[T1]{fontenc}
\usepackage{amsmath}
\usepackage{amsthm}
\usepackage{graphicx}
\usepackage{booktabs}
\usepackage{color}
\usepackage{xcolor}
\usepackage{comment}
\usepackage{subcaption}
\usepackage{caption}
\usepackage{algorithm}
\usepackage{algorithmic}
\usepackage{framed}
% \usepackage[linewidth=1pt]{mdframed}
\usepackage[top=35mm,headsep=15mm]{geometry}
\usepackage{ragged2e}
\justifying\let\raggedright\justifying
\usepackage{fancybox}
% \setlength{\parskip}{1mm}
% \setlength{\spaceskip}{\fontdimen2\font}

% \newdefinition{definition}{Definition}
% \newproof{proof}{Proof}

% \title{Difference Letter}

\pagestyle{empty}


\begin{document}
\centering
{\Large \textbf{Cover Letter}}

\raggedright
% \large
% \normalsize

% \vspace{5mm}

\noindent Dear Editor/Reviewer,  
% \vspace{1mm}
% \;\;\;\;\;\;\;\;\;\;\;\;\;\;\;\;\;\;\;\;\;\;\;\;\;\;\;\;\;\;\;\;\;\;\;\;\;\;\;\;\;\;\;\;\;\;\;\;\;\;\;\;\;\;\;\;\;\;\;\;\;\;\;\;\;\;\;\;\;\;\;\textcolor{red}{5 July 2023}{} 
%\\

We are pleased to submit the manuscript entitled, ``A Spatio-Temporal Adaptive Framework for Dispatching, Routing, and Pricing of Bus-Booking'', written by Yucen Gao, Pengfei Du, Xiaofeng Gao, Bin Wang and Xiaochun Yang, for consideration by your esteemed journal \emph{IEEE Transactions on Knowledge and Data Engineering (TKDE)}. We believe this work is well-suited for the scope of TKDE because it studies data-driven knowledge discovery and decision optimization for Intelligent Transportation Systems (ITS), with particular emphasis on spatio-temporal data modeling, robust learning under non-stationary demand, and joint optimization over real-world trajectory and order data.

Specifically, we propose ST-SACA, a novel end-to-end adaptive framework designed to solve the Joint Dispatching, Routing, and Pricing (JDRP) problem for bus-booking platforms. Unlike traditional methods that rely on stationary assumptions or treat operational decisions separately, ST-SACA learns robust policies from stochastic and dynamically changing demand environments. By incorporating a Spatio-Temporal Adaptive Actor, our model explicitly integrates geographical topology with demand dynamics through self-attention, enabling data-driven decision intelligence for coupled pricing, dispatching, and routing. Extensive experiments on real-world datasets demonstrate that ST-SACA significantly outperforms state-of-the-art baselines in terms of long-term profit and service reliability.

A preliminary version of this paper was presented at DASFAA 2023. This journal submission represents significant extensions: (1) Problem Reformulation: a generalized formulation distinguishing between Stationary and Robust JDRP to model stochastic uncertainty. (2) Spatio-Temporal Modeling Innovation: a redesigned architecture featuring the Spatio-Temporal Adaptive Actor to capture complex topological dependencies and temporal demand dynamics. (3) Theoretical Depth: new rigorous analyses regarding computational complexity, convergence properties, and robustness guarantees. (4) Extensive Evaluation: comprehensive new experiments evaluating the performance against state-of-the-art baselines under non-stationary scenarios. The specific differences between conference and journal versions will be emphasized clearly in the ``Difference Letter'' later.

% 本次期刊提交与会议版本之间存在显著差异，包括差异现有框架的不公平性、一些具体数据和示例、SACA的现实场景举例、JDRP问题的说明、目的站的数量选择以及$\lambda$对\textit{Rev}和\textit{ORR}的影响。



Correspondence should be addressed to the corresponding author, Prof. Xiaofeng Gao, at the following address, phone, and email address:
\vspace{2mm}

\doublebox{\parbox{130mm}{
\textbf{Xiaofeng Gao}, \vspace{0.5mm}\\ 
Telecom Building 3-407, Shanghai Jiao Tong University\\
800 Dongchuan Road, Minhang District, Shanghai, China, 200240\\ 
Phone:+86-21-3420-7407, \\
Email: gao-xf@cs.sjtu.edu.cn
}}
\vspace{2mm}

We appreciate your time and effort in reviewing this manuscript and look forward to receiving any helpful comments! Please feel free to contact us for any further information.

\noindent --------------

\noindent Sincerely yours,

\noindent Professor Xiaofeng Gao, PH.D.,

\noindent School of Computer Science

\noindent Shanghai Jiao Tong University, China
\end{document}
